\subsection{Puente cinemático--mecánico canónico}
\label{app:sp2-canonical-bridge}

Sea $\xi_L=(v_L^{\mathsf T},\omega)^{\mathsf T}$, sea
$J=\bigl[\begin{smallmatrix}0&-1\\1&0\end{smallmatrix}\bigr]$ y sea $r_i$ el
pivote de $i$ en el marco de la carga. Con radio $r_i^w>0$, semivía
$\ell_i^w>0$ y límite $\bar\omega_{w,i}$, la cinemática rígida exige
\begin{equation}
 v_i^{\mathrm{piv}}=v_L+\omega JR(\theta_L)r_i,
 \qquad
 \dot\varphi_{i,L/R}=\frac{v_{i,\parallel}\mp\ell_i^w\dot\theta_i}{r_i^w}.
 \label{eq:pre-sp2-pivot-wheels}
\end{equation}

\begin{proposicion}[Equivalencia cinemática]
\label{prop:pre-sp2-kinematic-equivalence}
En un intervalo con $v_i^{\mathrm{piv}}\ne0$, $\xi_L$ continuamente
diferenciable, pivotes inicialmente anclados y una rama continua de rumbo, la
coalición reproduce $\xi_L$ si y solo si cada eje longitudinal es paralelo o
antiparalelo a $v_i^{\mathrm{piv}}$, la velocidad lateral es nula y
\eqref{eq:pre-sp2-pivot-wheels} respeta los límites de rueda.
\end{proposicion}
\begin{proof}
Derivar $p_i=p_L+R(\theta_L)r_i$ da la primera igualdad; la restricción no
holónoma y el diferencial de ruedas dan la segunda. Recíprocamente, ambas
igualdades realizan la derivada del anclaje y la coincidencia inicial conserva
la relación rígida.
\end{proof}

En giro puro, $v_L=0$, la rueda exterior impone
$|\omega|\leq\min_{i\in\mathcal C}
r_i^w\bar\omega_{w,i}/(d_i^{\mathrm{ICR}}+\ell_i^w)$; el pivote situado en el
centro instantáneo se trata como el caso separado de giro sobre sí mismo.
Defínanse $\Xi_{\mathcal C}^{\mathrm{align}}=\bigcap_{i\in\mathcal C}
\Xi_i^{\mathrm{align}}$ y $\mathcal W_{\mathcal C}$ como el conjunto de
wrenches realizables por reacciones admisibles.

\begin{proposicion}[Monotonicidades opuestas]
\label{prop:pre-sp2-opposite-monotonicity}
Si $j\in\mathcal C$ y la reacción nula es admisible, entonces
$\Xi_{\mathcal C\setminus\{j\}}^{\mathrm{align}}\supseteq
\Xi_{\mathcal C}^{\mathrm{align}}$ y
$\mathcal W_{\mathcal C\setminus\{j\}}\subseteq\mathcal W_{\mathcal C}$.
La factibilidad conjunta no es, en general, monótona con la cardinalidad.
\end{proposicion}
\begin{proof}
Retirar $j$ elimina un conjunto de una intersección, pero también elimina sus
reacciones; la inclusión mecánica sigue al prolongar cualquier reparto restante
con reacción nula. Cada dirección puede ser estricta.
\end{proof}

\begin{proposicion}[Insuficiencia de la capacidad escalar]
\label{prop:pre-sp2-scalar-insufficiency}
Si $\|t_i\|\leq\bar t_i$ y $\sum_i t_i=F_d$, entonces
$\sum_i\bar t_i\geq\|F_d\|$ es necesario, pero no suficiente para realizar
$(F_d,\tau_d)$ ni para certificar soporte o fricción.
\end{proposicion}
\begin{proof}
La necesidad es triangular. Si todos los contactos están en el centro de masa,
$r_i\times t_i=0$ y ningún reparto genera $\tau_d\ne0$, por grande que sea la
suma escalar; omitir soporte o fricción solo relaja la factibilidad.
\end{proof}

El Teorema~\ref{prop:sp4-pose-stability} completa el enlace: con contactos
bilaterales fijos y wrench exacto, el equilibrio nominal de pose es localmente
asintóticamente estable. El resultado no cubre saturación persistente, pérdida
de contacto ni hardware.
